%%%%%%%%%%%%%%%%%%%%%%%%%%%%%%%%%%%%%%%%%%%%%%%%%%%%%%%%%%%%%%%%%%%%%%%%%%%%%%%%%%%%
\section{Domain of Values for Types\label{sec:DomainOfValuesForTypes}}
%%%%%%%%%%%%%%%%%%%%%%%%%%%%%%%%%%%%%%%%%%%%%%%%%%%%%%%%%%%%%%%%%%%%%%%%%%%%%%%%%%%%
This section formalises the concept of the set of values for a given type.
The formalism is given in the form of inference rules, although those are not meant
to be implemented.

We define the concept of a \emph{dynamic domain} of a type
and the \emph{static domain} of a type.
Intuitively, domains assign potentially infinite sets of \nativevaluesterm{} to types.
Dynamic domains are used by the semantics to evaluate expressions of the form \texttt{ARBITRARY: t}
by choosing a single value from the dynamic domain of $\vt$ (\SemanticsRuleRef{EArbitrary}).
Static domains are used to define subtype satisfaction (\TypingRuleRef{SubtypeSatisfaction})
via a conservative subsumption test as defined in \chapref{SymbolicDomainSubsetTesting}.

\subsection{The Dynamic Domain of a Type\label{sec:DynDomain}}
\RenderRelation{dynamic_domain}
%
We say that $\dynamicdomain(\env, \vt)$ is the \emph{dynamic domain} of $\vt$
in the environment $\env$.
%
The \emph{static domain} of a type is the set of values which storage elements of that type may hold
\underline{across all possible dynamic environments}.
%
The reason for this distinction is that the sets of values
of integer types, bitvector types, and array types can depend on the dynamic values of variables.

Types that do not refer to variables whose values are only known dynamically have
a static domain that is equal to any of their dynamic domains.
In those cases, we simply refer to their \emph{domain}.

Associating a set of values to a type is done by evaluating any expression appearing
in the type definitions.
%
Expressions appearing in types are guaranteed to be \sideeffectfreeterm{} by the
function $\annotatetype$.
%
Evaluation is defined by the relation $\evalexprsef\empty$,
which evaluates \sideeffectfreeexpressionsterm{} and either returns
a configuration of the form $\ResultExprSEF(\vv,\vg)$ or a dynamic error configuration $\DynErrorConfig$,
or diverges.
In the first case, $\vv$ is a \nativevalueterm{} and $\vg$
is an \emph{execution graph}. Execution graphs are related to the concurrent semantics
and can be ignored in the context of defining dynamic domains.
In the case of a \dynamicerrorterm{} (which can occur if, for example, an expression attempts to divide
\texttt{8} by \texttt{0}), a \DynamicErrorConfigurationTerm{}, for which we use the notation
$\DynErrorConfig$, is returned.
%
The dynamic domain is empty in cases where evaluating \sideeffectfreeexpressionsterm{}
results in a \DynamicErrorConfigurationTerm{} or diverges.
%
The dynamic domain is undefined if the type $\vt$ is not well-typed in $\tenv$.
That is, if $\annotatetype(\tenv, \vt) \typearrow \TypeErrorConfig$.

We use the following notations to define sets of values used as domains:
\hypertarget{relation-bitvectordomain}{}
\hypertarget{relation-vectordomain}{}
\hypertarget{relation-arraydomain}{}
\hypertarget{relation-recorddomain}{}
\[
\begin{array}{@{}ll@{}}
  \text{Expression} & \text{Set denoted by the expression}\\
  \hline
  \tbool &
    \{ \nvboolop{\True}, \nvboolop{\False} \}\\
  \tint &
    \{ \nvintop{n} \;|\; n \in \Z \}\\
  \treal &
    \{ \NVLiteral(\LReal(q)) \;|\; q \in \Q \}\\
  \tstring &
    \{ \nvstringop{s} \;|\; s \in \Strings \}\\
  \bitvectordomainop{k} &
    \begin{cases}
      \{ \nvbitvectorop{\vb_{1..k}} \;|\; \forall i=1..k: \vb_i \in \Bit \} & \text{if }k > 0\\
      \{ \nvbitvectorop{\emptylist} \} & \text{otherwise}
    \end{cases}\\
  \vectordomainop{D_{1..k}} &
    \{ \NVVector(\vv_{1..k}) \;|\; \vi=1..k: \vv_\vi \in D_\vi \}\\
  \arraydomainop[H]{k}{D} &
    \{ \NVVector(\vv_{1..k}) \;|\; \vi=1..k: \vv_\vi \in D \}\\
  \recorddomainop[H]{\id_{1..k}}{D_{1..k}} &
    \{ \NVRecord(\{ \vi=1..k: \id_\vi \mapsto \vv_\vi \}) \;|\; \vi=1..k: \vv_\vi \in D_\vi \}
\end{array}
\]

\RenderProseAndFormally{dynamic_domain}

\ExampleDef{Type Domains}
The domain of \texttt{integer} is the infinite set of all integers.

The domain of \verb|integer {2,16}| is the set $\{\nvint(2), \nvint(16)\}$.

The domain of \verb|integer{1..3}| is the set $\{\nvint(1), \nvint(2), \nvint(3)\}$.

The domain of \verb|integer{10..1}| is the empty set as there are no integers that are
both greater than $10$ and smaller than $1$.

The domain of \texttt{bits(2)} is the following set:
\[
  \left\{
  \begin{array}{l}
    \nvbitvector(00),\\
    \nvbitvector(01),\\
    \nvbitvector(10),\\
    \nvbitvector(11)
  \end{array}
  \right\}\enspace.
\]

The domain of \verb|enumeration {GREEN, ORANGE, RED}| is the set \\
$\{\nvlabel(\texttt{GREEN}), \nvlabel(\texttt{ORANGE}), \nvlabel(\texttt{RED})\}$ and so is the domain
of \\
\verb|type TrafficLights of enumeration {GREEN, ORANGE, RED}|.

The domain of \texttt{(integer, integer)} is the set containing all pairs of native integer values.

The domain of \verb|record {a: integer;  b: boolean}| contains all native records
that map \texttt{a} to a native integer value and \texttt{b} to a native Boolean value.

The dynamic domain of a subprogram parameter \texttt{N: integer} is the (singleton) set containing
the native integer value $c$,
which is assigned to \texttt{N} by a given dynamic environment. The static domain of that parameter
is the infinite set of all native integer values.

The domain of \texttt{array [[2]] of boolean} is the set
\[
  \left\{
  \begin{array}{l}
    \NVVector([\nvboolop{\False}, \nvboolop{\False}]),\\
    \NVVector([\nvboolop{\False}, \nvboolop{\True}]),\\
    \NVVector([\nvboolop{\True}, \nvboolop{\False}]),\\
    \NVVector([\nvboolop{\True}, \nvboolop{\True}])
  \end{array}
  \right\}.
\]

\subsection{The Dynamic Domain of a Constraint\label{sec:DynDomainConstraint}}
We also associate dynamic domains to integer constraints.
\RenderRelation{dynamic_domain_constraint}
\RenderProseAndFormally{dynamic_domain_constraint}

\ExampleDef{Constraint Domains}
The dynamic domain of the integer constraint \texttt{4} is the singleton set $\{\nvintop{4}\}$.
%
The dynamic domain of the integer constraint \texttt{1..3} is the set
$\{\nvintop{1}, \nvintop{2}, \nvintop{3}\}$.
